% Section 6: what the benchmark can and cannot test. Evidence: VANET PAPER2_SCOPE_REVIEW.md §4; protocol §8.8 (network fact,
% finding 4), §8.10 finding 5; ESWA plan §2 (NGSIM speed 6.24 m/s).
\section{What This Benchmark Can and Cannot Test}
\label{sec:testability}

The results above are statements about a method. This section is a statement about the benchmark the method was evaluated on, because the submitted paper's title-level claims---diverse, human-like and safe cooperative driving---were made on quantities this benchmark does not measure, and a reader who wants to reuse it, or one like it, should know which is which. Table~\ref{tab:t7_testability} gives the verdicts; three facts about the scenario explain most of them.

\paragraph{The exit task has a one-line solution.}
Of the four lanes of the first highway segment only lane~0 connects to the exit ramp; all four continue to the next segment. An exit-assigned agent therefore succeeds if and only if it is in lane~0 when the segment ends, and the policy ``command Lane\_Right at every step'' is optimal for it: 100\,\% exits at 24.5\,m/s with the lowest proximity of any policy we evaluated under mode 512 (Table~\ref{tab:t3_gate}). A learned policy cannot do better on exits than this constant, and a learned policy that does as well may simply \emph{be} this constant, as greedy anchored CALM is. Exit success is thus a one-bit test of whether a policy commands rightward moves. It is a useful bit---it separated the collapsed and oscillating policies from the others once lane changes were the policy's own---but it carries no information about gap selection, timing or yielding, the things a cooperative exit manoeuvre consists of.

\paragraph{Nothing in the task requires the agents to interact.}
Three agents have an exit to reach and two are mainline vehicles with no objective beyond the shared reward's speed and spacing terms. The agents depart one second apart into free-flowing traffic, each exit agent's best action does not depend on what the others do, and no reported metric measures an interaction: none counts yielding, gap creation or conflicts resolved. Five independent constant policies achieve the best exit, speed and proximity figures in this paper. A benchmark on which that is true cannot support a claim about cooperation, whatever the training algorithm shares among agents, and ``multi-agent'' describes how the policies were trained, not something the evaluation saw.

\paragraph{The human reference is a different traffic state.}
The NGSIM recording is a congested bottleneck with a mean speed of 6.24\,m/s in which a vehicle holds its lane and speed in more than half of the samples and changes lane in one sample in four hundred. The simulated task is free flow at 24--27\,m/s for a policy that does not obstruct itself, with a mandatory lane change for three of the five agents within one segment. Scoring a simulator policy by how often its action equals the recorded human action in the recorded human state therefore asks whether the policy would behave like a driver in a jam; the best possible answer under that score is to do nothing (55.3\,\%), which strands the exit agents in the simulator (7\,\% exits under mode 512). Agreement and exit success pull in opposite directions by construction. What can be compared with the recording is distributional: the per-state entropy of the policy on human states, and the divergence of its action marginal from the human one, both reported in Table~\ref{tab:t4_val}. Neither is a measure of human-likeness of behaviour, and we do not use the word for any policy in this paper.

\paragraph{What follows.}
The benchmark can test the mechanism the method is built on---whether an anchor keeps a prior's entropy through fine-tuning, against a matched control and against the prior itself---and it can describe behaviour with exposure-normalised proximity, speed, lane transitions and lane occupancy beside reference policies. It cannot test human-likeness, safety or cooperation. The first limitation is the dataset's, the second the simulator's, the third the scenario's, and none is repaired by more seeds or a better algorithm. A scenario in which they could be tested would need an exit policy that is not a constant (several lanes reaching the ramp, exits at different positions, or traffic dense enough that a rightward move is sometimes refused), goals that conflict so that coordination has content and a metric to measure it, a safety measure that the simulator's driver model does not already guarantee, and a human reference recorded in, or replayed into, the traffic state the task poses.

% T7: hand-written (not produced by make_r2_tables.py). Source: VANET PAPER2_SCOPE_REVIEW.md §4; protocol §8.6–§8.11.
\begin{table*}[!t]
\centering
\footnotesize
\caption{What the MultiExit\,+\,NGSIM benchmark can and cannot measure, with the evidence in this paper for each verdict.}
\label{tab:t7_testability}
\setlength{\tabcolsep}{4pt}
\begin{tabular}{@{}>{\raggedright\arraybackslash}p{0.19\textwidth}>{\raggedright\arraybackslash}p{0.13\textwidth}>{\raggedright\arraybackslash}p{0.49\textwidth}>{\raggedright\arraybackslash}p{0.13\textwidth}@{}}
\toprule
Quantity & Measurable here? & Why & Evidence \\
\midrule
Per-state entropy on held-out human states (Definition~\ref{def:collapse}) & Yes & A dataset-level quantity, independent of the simulator; computed per actor from the softmax output on a vehicle-grouped validation split under a recorded contract. & \S\ref{sec:agreement}, \S\ref{sec:anchor} \\
Per-state entropy over simulator rollouts & Yes, under a declared execution rule & Greedy rollouts of the anchored policy visit the constant policy's states (0.64--0.66); sampled rollouts give 0.71. The rule changes the states visited and must be fixed in advance. & \S\ref{sec:concealed}, \S\ref{sec:anchor} \\
Resistance to collapse against a matched control & Yes & Fixed initialisation, one lane-change mode for training and evaluation, one selection rule per arm, matched episode seeds, paired bootstrap, a decision rule written first. & \S\ref{sec:anchor} \\
Improvement over the frozen prior & Yes & The prior is a policy and can be run through the same evaluator on the same episodes. & \S\ref{sec:prior} \\
Proximity exposure per agent-km; speed & As behaviour descriptors only & They track speed and lane behaviour (sampled policies three to five times greedy ones, oscillators ten times, stranded policies a hundred times) and are meaningful only beside the constant-action and SUMO-default rows; they are not safety outcomes. & \S\ref{sec:concealed} \\
Exit success & As a one-bit test & Only lane~0 reaches the ramp, so ``always move right'' is optimal for an exit agent; under mode 512 exit success measures whether a policy commands rightward moves, not how well agents coordinate. Under SUMO's default mode it measures SUMO. & \S\ref{sec:concealed} \\
Human-likeness of behaviour & No & Agreement is a three-class problem with a 55.3\,\% trivial baseline that no method reaches; lane changes are 0.26\,\% of the labels; the recording is congested traffic near 6\,m/s and the simulated task is free-flow, so there is no behavioural human reference inside the simulator. & \S\ref{sec:agreement} \\
Collisions & No & SUMO's car-following and gap enforcement prevent them for every policy, stranded and oscillating ones included: none in roughly 4{,}100 episodes. & \S\ref{sec:concealed} \\
Cooperation & No & No metric measures interaction between agents; two of the five agents have no task; nothing distinguishes a coordinated solution from five independent constant policies. & \S\ref{sec:testability} \\
Reward & Training statistic only & The evaluator returns none; last-episode training means under exploration are comparable across arms trained under one contract and are not deployment performance. & \S\ref{sec:reward} \\
\bottomrule
\end{tabular}
\end{table*}

